\documentclass[11pt]{article}
\usepackage[T1]{fontenc}
\usepackage[utf8]{inputenc}
\usepackage{amsmath,amssymb,amsthm,mathtools}
\usepackage{booktabs}
\usepackage{geometry}
\usepackage{microtype}
\usepackage[hidelinks]{hyperref}
\geometry{margin=1in}

\newtheorem{theorem}{Theorem}
\newtheorem{proposition}[theorem]{Proposition}
\newtheorem{corollary}[theorem]{Corollary}
\newtheorem{lemma}[theorem]{Lemma}
\theoremstyle{definition}
\newtheorem{definition}[theorem]{Definition}
\newtheorem{remark}[theorem]{Remark}

\newcommand{\PK}{P_K}
\newcommand{\QK}{Q_K}
\newcommand{\eps}{\varepsilon}
\newcommand{\Torus}{\mathbb T^3}

\hypersetup{
  pdftitle={Finite Readout Certificates for Navier--Stokes},
  pdfauthor={Yaoharee Lahtee},
  pdfsubject={Discrete epsilon-completion; Navier--Stokes; a posteriori certification},
  pdfkeywords={Navier--Stokes, Fourier-Galerkin, epsilon-completion, Leray-Hopf, spectral tail, relative energy, a posteriori error}
}

\title{Finite Readout Certificates for Navier--Stokes:\\
Discrete Epsilon-Completion, Spectral Tails, and A Posteriori Adapters}
\author{Yaoharee Lahtee\\
\small Open Civil Science Initiative, Bangkok, Thailand}
\date{September 2026}

\begin{document}
\maketitle

\begin{abstract}
Finite numerical agreement is not, by itself, a certificate that an infinite-dimensional
Navier--Stokes state has been reconstructed to a declared tolerance.  This paper develops a
fail-closed formulation of that distinction.  A nested finite-refinement defect
$\delta_K$ measures disagreement on shared retained information, while a separate quantity
$\beta_{K,Y}$ is required to bound distinctions omitted by the finite representation in a
declared target space $Y$.  We first prove an elementary obstruction: terminal Fourier
coefficients below a cutoff cannot, over an unrestricted divergence-free $L^2$ class, determine a
universal useful bound on the omitted terminal tail.  We then identify target spaces and retained
records for which rigorous bounds do exist.  For an unforced Leray--Hopf solution on the periodic
three-torus,
\[
\|(I-P_K)u\|_{L^2(0,T;L^2_x)}
\le \frac{\|u_0\|_2}{\sqrt{2\nu}\,(K+1)},
\]
which yields a computable $O(K^{-1})$ spacetime certificate without assuming global smoothness.
A terminal $L^2$ certificate follows either from a separately certified pointwise $H^s$ bound or
from a richer retained energy/dissipation tape.  Finally, a residual-based relative-energy
estimate connects a smooth finite Fourier comparison path to a Leray--Hopf solution and isolates
the remaining numerical obligation: a validated continuous-time enclosure of the actual RK4
tape.  The associated finite Fourier residual is itself a finite triad computation.  Executable
finite diagnostics and a cross-repository provenance ledger accompany the analysis.  None of the
results proves global regularity, finite-time blow-up, uniqueness of weak solutions, or any part of
the Clay Millennium Navier--Stokes problem.
\end{abstract}

\section{Introduction}

A finite computation answers a finite question.  A continuum claim requires an additional bridge.
This elementary observation becomes operationally important in numerical PDE work, where several
distinct statements are easily conflated:
\[
\text{finite algebra correct}
\quad\not\Rightarrow\quad
\text{continuum complete}
\quad\not\Rightarrow\quad
\text{physically validated turbulence}.
\]
The present paper studies how the middle implication can be made explicit and fail-closed for
incompressible Navier--Stokes computations.

The starting point is the Discrete Epsilon-Completion (EPSC) programme.  At refinement level $K$,
a finite state or readout $x_K$ is compared with the restriction of the next refinement.  Small
nested disagreement is useful evidence, but it does not control distinctions that both finite
representations omit.  We therefore separate two quantities:
\begin{align}
\delta_K &= \|R_K x_{K+1}-x_K\|, \\
\beta_{K,Y} &= \text{a proved upper bound on omitted information in target }Y.
\end{align}
Only when the relevant hypotheses are certified does the fail-closed gate permit
\[
\delta_K+\beta_{K,Y}\le\eps
\quad\Longrightarrow\quad
\mathrm{CERTIFIED}_\eps(Y).
\]
If the required $\beta_{K,Y}$ is unavailable, the verdict is \texttt{HOLD} regardless of how small
$\delta_K$ appears numerically.

The contribution here is not a new energy inequality or a new weak--strong stability theorem.
Those analytic ingredients are classical.  The contribution is the explicit separation of target,
retained record, adapter, and omitted-information certificate, together with a concrete Fourier
implementation path and machine-executable finite checks.  This separation identifies precisely
which parts of the finite-to-continuum claim are already certified and which remain open.

\section{Setting and notation}

Let $u$ solve incompressible Navier--Stokes on the $2\pi$-periodic three-torus $\Torus$,
\begin{equation}\label{eq:NS}
\partial_t u + \mathbb P[(u\cdot\nabla)u] - \nu\Delta u = \mathbb P f,
\qquad \nabla\cdot u=0,
\end{equation}
with viscosity $\nu>0$, where $\mathbb P$ is the Leray projector.  We use the Fourier
normalization
\[
\|u\|_2^2=\sum_{k\in\mathbb Z^3}|\widehat u_k|^2,
\qquad
\|\nabla u\|_2^2=\sum_{k\in\mathbb Z^3}|k|^2|\widehat u_k|^2.
\]
Let $P_K$ retain the cube $\|k\|_\infty\le K$ and write $Q_K=I-P_K$.  Every omitted mode then
satisfies $|k|_2\ge K+1$.

For the finite computation, the production representation is a set of integer mode labels and
complex divergence-free coefficients.  A finite Fourier--Galerkin trajectory obeys
\begin{equation}\label{eq:galerkin}
\frac{d\widehat v_k}{dt}
=-\nu|k|^2\widehat v_k
-iP_k\sum_{p+q=k}(q\cdot\widehat v_p)\widehat v_q,
\qquad \|k\|_\infty\le K,
\end{equation}
with all sums finite after the cutoff is fixed.  The implementation used in the accompanying
repository evaluates these ordered triads directly and does not require a spatial grid in the
production path.  This is a computational statement about the finite recurrence, not an
ontological statement that continuum space has been eliminated.

\section{Discrete epsilon-completion}

\begin{definition}[Nested consistency defect]
Given a restriction map $R_K$ from refinement $K+1$ to the information represented at level $K$,
define
\begin{equation}\label{eq:delta}
\delta_K(x):=\|R_Kx_{K+1}-x_K\|.
\end{equation}
\end{definition}

For Fourier Navier--Stokes diagnostics, one may also record the outer retained-shell energy
\begin{equation}\label{eq:boundary}
E_{\partial K}(u)=\frac12\sum_{\|k\|_\infty=K}|\widehat u_k|^2.
\end{equation}
This is a finite diagnostic proxy for unresolved activity.  It is not a theorem-level substitute
for a bound on the full omitted tail.

\begin{definition}[Target-indexed omitted-information certificate]
Let $Y$ be a declared target norm or readout space and $\mathcal R_K$ a declared finite retained
record under assumptions $\mathcal A$.  An EPSC tail certificate is a computable quantity
$\beta_{K,Y}(\mathcal R_K;\mathcal A)\ge0$ such that
\begin{equation}\label{eq:beta}
\|Q_Kx\|_Y\le \beta_{K,Y}(\mathcal R_K;\mathcal A).
\end{equation}
For asymptotic completion one additionally requires $\beta_{K,Y}\to0$ along the declared
refinement sequence.
\end{definition}

This definition makes explicit that ``complete'' is incomplete language unless the reader/norm,
admissible class, and available retained record are specified.

\section{Terminal finite-record obstruction}

\begin{proposition}[Low terminal modes do not identify an unrestricted tail]\label{prop:nogo}
Fix $K\ge0$.  Over the class of divergence-free, real-valued $L^2$ Fourier states, the retained
terminal record $P_Ku$ alone does not determine any finite universal upper bound on
$\|Q_Ku\|_2$.
\end{proposition}

\begin{proof}
Choose $q=(K+1,0,0)$ and a nonzero vector $a\perp q$.  Define a real divergence-free perturbation
$v$ by placing Fourier coefficients only at $q$ and $-q$, with
$\widehat v_q=a$ and $\widehat v_{-q}=\overline a$.  Then $P_Kv=0$, while
$\|Q_Kv\|_2=\sqrt{2}|a|$.  Scaling $a$ leaves the retained record unchanged and makes the omitted
tail arbitrarily large.  Hence no function of $P_Ku$ alone can give a finite universal terminal
bound on the unrestricted class.
\end{proof}

\begin{remark}
Proposition~\ref{prop:nogo} is an identifiability statement about a terminal state class.  It does
not assert that Navier--Stokes dynamics realizes every such perturbation at a prescribed time.
Rather, it shows why a useful terminal certificate must contain additional admissible-class or
dynamical information.
\end{remark}

\section{A positive spacetime certificate}

\begin{lemma}[Spectral tail inequality]\label{lem:spectral}
For $s>0$ and any $u$ with finite homogeneous $H^s$ seminorm,
\begin{equation}\label{eq:Hs}
\|Q_Ku\|_2
\le \frac{\||\nabla|^s u\|_2}{(K+1)^s}.
\end{equation}
\end{lemma}

\begin{proof}
Every omitted mode has $|k|\ge K+1$.  Therefore
\[
\sum_{\|k\|_\infty>K}|\widehat u_k|^2
\le (K+1)^{-2s}
\sum_{\|k\|_\infty>K}|k|^{2s}|\widehat u_k|^2,
\]
and taking square roots gives the result.
\end{proof}

\begin{theorem}[Unforced Leray--Hopf spacetime tail certificate]\label{thm:spacetime}
Let $u$ be a Leray--Hopf weak solution of the unforced equation on $[0,T]$.  Then
\begin{equation}\label{eq:spacetime}
\boxed{
\|Q_Ku\|_{L^2(0,T;L^2_x)}
\le
\beta_K^{\mathrm{ST}}
:=\frac{\|u_0\|_2}{\sqrt{2\nu}\,(K+1)}
}
\end{equation}
and $\beta_K^{\mathrm{ST}}\to0$ as $K\to\infty$.
\end{theorem}

\begin{proof}
Lemma~\ref{lem:spectral} with $s=1$ gives
\[
\|Q_Ku(t)\|_2^2\le (K+1)^{-2}\|\nabla u(t)\|_2^2.
\]
The Leray--Hopf energy inequality gives
\[
\frac12\|u(t)\|_2^2
+\nu\int_0^t\|\nabla u(s)\|_2^2\,ds
\le\frac12\|u_0\|_2^2.
\]
Integrating the spectral bound and discarding the nonnegative terminal kinetic energy yields
\[
\int_0^T\|Q_Ku\|_2^2dt
\le \frac{\|u_0\|_2^2}{2\nu(K+1)^2}.
\]
Taking square roots proves~\eqref{eq:spacetime}.
\end{proof}

For forcing $f\in L^2(0,T;H^{-1})$, Young's inequality gives the conservative variant
\begin{equation}\label{eq:forced}
\|Q_Ku\|_{L^2_tL^2_x}
\le
\frac{1}{K+1}
\left(
\frac{\|u_0\|_2^2}{\nu}
+\frac{\|f\|_{L^2_tH^{-1}}^2}{\nu^2}
\right)^{1/2}.
\end{equation}

\begin{corollary}[Lipschitz readout lift]\label{cor:lipschitz}
If a declared readout $\mathcal Q:Y\to Z$ is $L_{\mathcal Q}$-Lipschitz in a space $Y$ for which an
EPSC certificate $\|u-P_Ku\|_Y\le\beta_{K,Y}$ has been established, then
\[
\|\mathcal Q(u)-\mathcal Q(P_Ku)\|_Z
\le L_{\mathcal Q}\beta_{K,Y}.
\]
\end{corollary}

For example, for the time average $\bar u=T^{-1}\int_0^T u(t)dt$,
Cauchy--Schwarz gives
\[
\|Q_K\bar u\|_2\le \beta_K^{\mathrm{ST}}/\sqrt T.
\]

\section{Terminal certificates from richer retained records}

\subsection{A conditional $H^s$ certificate}

If a separate argument certifies
\[
\||\nabla|^su(T)\|_2\le M_s(T),
\]
then Lemma~\ref{lem:spectral} immediately gives
\begin{equation}\label{eq:terminalHs}
\boxed{
\|Q_Ku(T)\|_2\le\frac{M_s(T)}{(K+1)^s}.
}
\end{equation}
The difficult part in three dimensions is not the spectral inequality but obtaining an appropriate
globally valid pointwise regularity bound.  EPSC therefore does not bypass the classical regularity
obstacle for arbitrary terminal reconstruction.

\subsection{An energy/dissipation tape}

The terminal no-go concerns low terminal coefficients alone.  The energy inequality provides a
natural richer retained record.

\begin{theorem}[Terminal energy-budget certificate]\label{thm:energybudget}
Let $u$ be an unforced Leray--Hopf solution.  Suppose finite certified data satisfy
\[
U_0\ge\|u_0\|_2,
\qquad
L_K(T)\le\|P_Ku(T)\|_2,
\]
and
\[
D_K(T)\le
\nu\int_0^T\|\nabla P_Ku(t)\|_2^2dt.
\]
Then, whenever the right-hand side is nonnegative,
\begin{equation}\label{eq:energybudget}
\boxed{
\|Q_Ku(T)\|_2
\le
\beta_K^{\mathrm{EB}}
:=
\left[U_0^2-L_K(T)^2-2D_K(T)\right]^{1/2}.
}
\end{equation}
A materially negative bracket signals inconsistent directional certificates and must produce a
fail-closed verdict rather than being clamped to zero.
\end{theorem}

\begin{proof}
Orthogonality of Fourier projection gives
\[
\|u(T)\|_2^2=\|P_Ku(T)\|_2^2+\|Q_Ku(T)\|_2^2,
\]
and likewise
\[
\|\nabla u\|_2^2
=\|\nabla P_Ku\|_2^2+\|\nabla Q_Ku\|_2^2.
\]
Insert both decompositions into the energy inequality and discard the nonnegative unresolved
dissipation term.  Replacing exact quantities by bounds in the indicated directions yields
~\eqref{eq:energybudget}.
\end{proof}

With exact projected quantities, define the energy-inequality slack
\begin{equation}\label{eq:defect}
\mathcal D_E(T)
:=\|u_0\|_2^2-\|u(T)\|_2^2
-2\nu\int_0^T\|\nabla u\|_2^2dt\ge0.
\end{equation}
Then the exact squared budget decomposes as
\begin{equation}\label{eq:floor}
(\beta_K^{\mathrm{EB}})^2
=\|Q_Ku(T)\|_2^2
+2\nu\int_0^T\|\nabla Q_Ku\|_2^2dt
+\mathcal D_E(T).
\end{equation}
Consequently
\[
\lim_{K\to\infty}(\beta_K^{\mathrm{EB}})^2=\mathcal D_E(T).
\]
If energy equality holds, $\mathcal D_E(T)=0$ and this particular terminal certificate tends to
zero.  For a general Leray--Hopf solution the actual Fourier tail still tends to zero at each fixed
$L^2$ time, but the energy-budget certificate can retain a nonzero floor because the energy
inequality alone does not account for the slack.

\section{A residual-based Galerkin-to-continuum adapter}

Theorem~\ref{thm:energybudget} concerns the actual continuum projection.  A numerical Galerkin
trajectory must not be silently substituted for $P_Ku$.  A standard relative-energy estimate gives
an explicit adapter.

Let $v$ be a smooth divergence-free comparison path and define its full-equation residual
\begin{equation}\label{eq:residual}
r
:=\partial_t v+\mathbb P[(v\cdot\nabla)v]-\nu\Delta v-\mathbb P f.
\end{equation}
Let $e_0\ge\|u_0-v(0)\|_2$ and define
\begin{equation}\label{eq:AB}
A_T:=2\int_0^T\|\nabla v(t)\|_\infty dt,
\qquad
B_T:=\int_0^T\|r(t)\|_{H^{-1}}^2dt.
\end{equation}

\begin{theorem}[Relative-energy adapter]\label{thm:adapter}
Let $u$ be a Leray--Hopf weak solution of~\eqref{eq:NS} and $v$ a sufficiently smooth,
divergence-free comparison path with the same boundary conditions and forcing convention.  Then
\begin{equation}\label{eq:adapter}
\boxed{
\sup_{0\le t\le T}\|u(t)-v(t)\|_2^2
\le
e^{A_T}\left(e_0^2+\frac{B_T}{\nu}\right).
}
\end{equation}
\end{theorem}

\begin{proof}
Set $w=u-v$.  Testing the difference equation with $w$ and using the incompressible cancellation
leaves the deformation term $\langle (w\cdot\nabla)v,w\rangle$.  Therefore
\[
\frac12\frac{d}{dt}\|w\|_2^2+\nu\|\nabla w\|_2^2
\le \|\nabla v\|_\infty\|w\|_2^2
+\|r\|_{H^{-1}}\|\nabla w\|_2.
\]
Young's inequality absorbs half the viscous term and gives
\[
\frac{d}{dt}\|w\|_2^2
\le 2\|\nabla v\|_\infty\|w\|_2^2
+\nu^{-1}\|r\|_{H^{-1}}^2.
\]
Gronwall's inequality and the definitions~\eqref{eq:AB} prove~\eqref{eq:adapter}.
\end{proof}

If $v(T)$ is supported in $P_K$, then $Q_Kv(T)=0$, so
\begin{equation}\label{eq:relativebeta}
\boxed{
\|Q_Ku(T)\|_2
\le
\beta_K^{\mathrm{RE}}
:=e^{A_T/2}
\left(e_0^2+\frac{B_T}{\nu}\right)^{1/2}.
}
\end{equation}
This is a terminal continuum bound obtained from a comparison path and certified residual summary;
it does not require identifying the numerical state with the continuum projection.

\section{The unresolved residual is a finite triad tape}

For an exact $K$-supported Fourier--Galerkin path $v$, the retained equations cancel the projected
part of~\eqref{eq:residual}.  The nonlinear product can only generate modes in the finite
$2K$ cube.  For omitted output mode $q$,
\begin{equation}\label{eq:residualmode}
\widehat r_q
=
P_q\left[
 i\sum_{p+s=q}(s\cdot\widehat v_p)\widehat v_s
\right],
\qquad \|q\|_\infty>K,
\end{equation}
and
\begin{equation}\label{eq:hminusone}
\|r\|_{H^{-1}}^2
=
\sum_{q\ne0}\frac{|\widehat r_q|^2}{|q|^2}.
\end{equation}
For fixed $K$ these are finite sums.  Likewise, a conservative Fourier-polynomial bound is
\begin{equation}\label{eq:gradlinf}
\|\nabla v\|_\infty
\le
\sum_{k}|k|\,|\widehat v_k|.
\end{equation}
Thus the analytic adapter reduces the remaining end-to-end numerical work to certified time
integration/enclosure of finite quantities.

\section{Validated RK4 continuous-time enclosure: the remaining frontier}

The current finite solver uses explicit RK4 at discrete nodes.  Nodewise residual samples do not
certify the integrals in~\eqref{eq:AB}.  The remaining numerical obligation is therefore precise:
construct from the floating-point RK4 tape a validated continuous-time finite Fourier path
$v_h(t)$ and computable enclosures
\begin{equation}\label{eq:enclosure}
A_T\le\overline A_T,
\qquad
B_T\le\overline B_T.
\end{equation}
Once~\eqref{eq:enclosure} is available together with a certified initial mismatch, the
relative-energy theorem gives the end-to-end bound
\begin{equation}\label{eq:endtoend}
\|Q_Ku(T)\|_2
\le
e^{\overline A_T/2}
\left(e_0^2+\frac{\overline B_T}{\nu}\right)^{1/2}.
\end{equation}

A practical validated implementation can use a piecewise-polynomial continuous extension of each
RK4 step, outward-rounded interval bounds for its coefficients, interval evaluation of the finite
triad residual on each time cell, and certified quadrature/enclosure of the resulting nonnegative
polynomial or interval majorant.  This paper does not claim that this final interval-enclosure layer
has already been completed.  In the accompanying Toledo registry it remains the explicit open
target \texttt{PROP-EPSC-15}.

\section{Executable finite evidence}

The analytic claims above are deliberately separated from finite numerical diagnostics.  The
accompanying reproduction system currently records the following finite checks.

\begin{table}[h]
\centering
\small
\begin{tabular}{lll}
\toprule
Check & Recorded result & Tier \\
\midrule
Analytic shear finite recurrence & max error $3.886\times10^{-16}$ & finite diagnostic \\
Direct triad vs padded FFT, $K=2$ & RHS diff $1.041\times10^{-16}$ & finite diagnostic \\
One RK4 step, direct vs FFT & diff $8.674\times10^{-19}$ & finite diagnostic \\
Taylor--Green nested gate, $K=4$ & $\delta_K=3.551\times10^{-8}$ & finite diagnostic \\
Taylor--Green nested gate, $K=5$ & $\delta_K=7.741\times10^{-10}$ & finite diagnostic \\
Spectral tail algebra & PASS on deterministic finite spectra & finite diagnostic \\
Terminal no-go witness & PASS; high-mode amplitude scalable & finite diagnostic \\
Energy-budget finite algebra & PASS & finite diagnostic \\
Relative-energy finite helper algebra & PASS & finite diagnostic \\
\bottomrule
\end{tabular}
\caption{Finite executable evidence.  These checks validate finite algebra or recorded numerical
instances; they do not upgrade the analytic theorems to machine proofs and do not certify turbulent
DNS adequacy.}
\label{tab:evidence}
\end{table}

The short Taylor--Green nested-refinement run used $\nu=0.01$, $\Delta t=0.005$, and $T=0.05$.
The operational nested diagnostic passed at $K=4$ and was confirmed at $K=5$ under its declared
thresholds.  This does not imply that $K=5$ contains all continuum information.  A separate
external-reference test at coarse turbulent resolution previously found that $K\le3$ did not
reproduce the published $Re=1600$ Taylor--Green dissipation curve; the present short-horizon
stability check does not overturn that result.

\section{Repository architecture and provenance}

The research is intentionally split across four public repositories so that interpretation,
provenance, executable mathematics, and domain-specific evidence cannot silently collapse into one
another:
\begin{itemize}
\item \texttt{information-discrete-math}: reusable EPSC algorithms and certificate helpers;
\item \texttt{toledo}: proposal/equation lineage and verification status;
\item \texttt{readout-problem-navier-stokes}: Navier--Stokes analysis, manuscript, finite evidence,
      reproduction ledger, and claim boundary;
\item \texttt{readout\_genesis}: interpretation/application map only; NS results are not promoted to
      root ontology.
\end{itemize}
The EPSC lineage is registered as \texttt{PROP-EPSC-01} through \texttt{PROP-EPSC-15}.  Classical
Fourier, energy, Sobolev, Leray--Hopf and relative-energy inequalities are registered for lineage
and role, not claimed as original mathematics.

AI systems were used for drafting, code integration, adversarial consistency checking, and
reproducibility work.  The human author is responsible for the mathematical claims, source
selection, interpretation, and final release decisions.  AI assistance is not treated as authorship
or as certification of truth.

\section{Claim boundary}

The strongest positive statement in this paper is target-specific.  A rigorous omitted-tail
certificate exists for Leray--Hopf trajectories in $L^2_tL^2_x$, and conditional/a-posteriori
terminal certificates exist from declared richer information.  A residual-based adapter converts
a certified smooth comparison path into a terminal $L^2$ bound.  The current numerical solver is
not yet an end-to-end continuum certificate because the validated continuous-time enclosure of its
floating-point RK4 tape remains open.

In particular, this work does \emph{not} establish global smoothness, finite-time singularity,
uniqueness of Leray--Hopf weak solutions, an autonomous finite closure of three-dimensional
Navier--Stokes, universal computational speedup, or turbulent DNS accuracy at coarse cutoff.  It
does not solve or partially solve the Clay Millennium existence-and-smoothness problem.  Its role
is narrower: it turns an ambiguous statement that a finite calculation ``looks converged'' into a
set of explicit mathematical obligations that can either be certified or fail closed.

\section{Conclusion}

The finite-to-infinite gap in Navier--Stokes computation is not one question.  It depends on what is
being read, in which norm, from which retained record, under which assumptions, and through which
adapter.  Terminal low-mode coefficients alone are non-identifying on an unrestricted class.  By
contrast, the Leray--Hopf energy inequality yields an explicit $O(K^{-1})$ spacetime certificate;
a richer energy/dissipation tape yields a terminal budget bound; and a residual-based
relative-energy estimate turns a certified finite comparison path into a terminal continuum bound.
The remaining implementation frontier is now finite and concrete: validated continuous-time
residual enclosure of the RK4 tape.  This separation is the central practical result of the EPSC
programme.

\begin{thebibliography}{9}
\bibitem{Leray1934}
J. Leray,
\newblock Sur le mouvement d'un liquide visqueux emplissant l'espace,
\newblock \emph{Acta Mathematica} 63 (1934), 193--248.

\bibitem{Hopf1951}
E. Hopf,
\newblock "Uber die Anfangswertaufgabe f\"ur die hydrodynamischen Grundgleichungen,
\newblock \emph{Mathematische Nachrichten} 4 (1951), 213--231.

\bibitem{Temam}
R. Temam,
\newblock \emph{Navier--Stokes Equations: Theory and Numerical Analysis},
\newblock AMS Chelsea Publishing.

\bibitem{ConstantinFoias}
P. Constantin and C. Foias,
\newblock \emph{Navier--Stokes Equations},
\newblock University of Chicago Press, 1988.

\bibitem{Fefferman}
C. L. Fefferman,
\newblock Existence and smoothness of the Navier--Stokes equation,
\newblock Clay Mathematics Institute Millennium Problem description, 2000.

\bibitem{Lahtee2026}
Y. Lahtee,
\newblock A Readout Problem for Fefferman's Existence and Smoothness of the Navier--Stokes Equation,
\newblock Zenodo, DOI: 10.5281/zenodo.22673246, 2026.
\end{thebibliography}

\end{document}
